\chapter{The specular coefficient \texorpdfstring{$F_{\bm 0}$}{F0}}
\label{chap:specular}

\section{What is different at \texorpdfstring{$\bm h=\bm 0$}{h = 0}}

Everything that distinguishes the specular rod was already established in
Eqs.~\eqref{eq:coeff-nonspec}--\eqref{eq:coeff-spec}: the prepared reference
$n_0^2(z)$ is laterally uniform, so it appears in the $\bm h=\bm 0$ Fourier
coefficient of $\delta n^2$ and in no other.  Equivalently, in the language of
Eq.~\eqref{eq:lateral-delta}, a laterally uniform density has all of its
lateral Fourier weight at $\bm q-\bm k_{i\parallel}=\bm 0$, and a genuine rod
$\bm h\neq\bm 0$ never sits there.

Consequently $F_{\bm 0}$, and only $F_{\bm 0}$, receives a second contribution
besides the atomic sum of Eq.~\eqref{eq:Fh-atomic}:
\begin{equation}
  F_{\bm 0}=\sum_u\left(F_{\bm 0,u}^{\rm at}+F_{\bm 0,u}^{\rm ref}\right),
  \qquad
  F_{\bm h}=\sum_uF_{\bm h,u}^{\rm at}\ \ (\bm h\neq\bm 0).
  \label{eq:Fh-two-parts}
\end{equation}
This is one matrix element with two kinds of contribution, not two matrix
elements.  In particular nothing is ``subtracted afterwards'': the sign is
already carried by $\Delta\rho_f$ in Eq.~\eqref{eq:density-contrast}.

It is worth stating what would go wrong if the reference contribution were
instead applied cell by cell off-specular.  Doing so would replace the exact
lateral delta function of Eq.~\eqref{eq:lateral-delta} by a finite-cell shape
transform and manufacture diffuse scattering that the reference profile does
not produce.  Instrumental resolution may mix genuine diffuse intensity into a
detector pixel on the specular rod, but that belongs to a later resolution
convolution, not to an interpolation of Eq.~\eqref{eq:Fh-two-parts} near
$\bm h=\bm 0$.

\section{Integral of one constant optical slab}

Within prepared medium $m$, both $\overline\rho_{f,m}$ and every channel
coefficient are constant; only the propagation phase varies with $z$.  Define
the stable depth integral
\begin{equation}
  \Phiint(q;a,b)
  =\int_a^b\ee^{\ii qz}\,\dd z
  =\begin{cases}
    \dfrac{\ee^{\ii qb}-\ee^{\ii qa}}{\ii q}, & q\neq0,\\[3mm]
    b-a, & q=0.
  \end{cases}
  \label{eq:phi-finite}
\end{equation}
For numerical work, an equivalent midpoint form is preferable:
\begin{equation}
  \Phiint(q;a,b)
  =(b-a)\,
   \ee^{\ii q(a+b)/2}
   \operatorname{sinc}\!\left(\frac{q(b-a)}{2}\right),
  \qquad
  \operatorname{sinc}x=\frac{\sin x}{x},
  \label{eq:phi-sinc}
\end{equation}
or the numerator of Eq.~\eqref{eq:phi-finite} may be evaluated with a complex
\texttt{expm1}.  For a lower terminal half-space, with the physical decaying
branch,
\begin{equation}
  \Phiint(q;-\infty,b)=\frac{\ee^{\ii qb}}{\ii q}.
  \label{eq:phi-half-space}
\end{equation}

Carrying the $-\sum_m\overline\rho_{f,m}\mathbf 1_{I_m}$ term of
Eq.~\eqref{eq:density-contrast} through Eq.~\eqref{eq:Fh-definition}, the
reference contribution to $F_{\bm 0}$ at profile level is
\begin{equation}
  \boxed{
  \sum_uF_{\bm 0,u}^{\rm ref}
  =-A_{\rm ref}
   \sum_m\overline\rho_{f,m}
   \sum_{\nu\in\{++,+-,-+,--\}}
   \Cchan_{m,\nu}
   \Phiint\!\left(Q_{z,m}^{\nu};z_m^-,z_m^+\right).
  }
  \label{eq:profile-reference-term}
\end{equation}
Finite graded layers generally retain all four channels in
Eq.~\eqref{eq:profile-reference-term}; only the terminal substrate reduces to
$++$, by Eq.~\eqref{eq:terminal-branches}.

\section{A well-defined per-record partition}

Only the total reference profile is physical.  A decomposition of that profile
among overlapping model objects is bookkeeping and is not unique unless a
rule is fixed.  orGUI can make the decomposition deterministic by retaining
the forward-density contribution of every generated record while building the
profile.

Let $\overline\rho_{f,u,m}$ denote record $u$'s share of the constant reference
density in prepared medium $m$.  The required conservation rule is
\begin{equation}
  \boxed{
  \sum_u\overline\rho_{f,u,m}=\overline\rho_{f,m}
  \quad\text{for every prepared medium }m.
  }
  \label{eq:reference-partition}
\end{equation}
Before optical-profile simplification, a generated slab of physical thickness
$t_{ud}$ contributes naturally
\begin{equation}
  \overline\rho_{f,ud}
  =\frac{w_u p_{ud}}{A_u t_{ud}}
   \sum_{a\in u}o_a f_a(0,E).
  \label{eq:natural-density-share}
\end{equation}
Normal strain changes $t_{ud}$ and hence the volume density, but preserves the
integrated areal content.  Signed occupancies in interface and etching models
remain signed in Eq.~\eqref{eq:natural-density-share}.

Coincident profile contributions are added.  If adjacent optical media are
merged, first determine the merged boundaries from the total profile, then
thickness-average every record's density over those same boundaries.  This is
a linear operation once the grouping is fixed and preserves
Eq.~\eqref{eq:reference-partition}.  Reconstructing shares independently with
different merge decisions would not preserve the prepared reference.

With Eq.~\eqref{eq:reference-partition} fixed, record $u$'s reference
contribution is
\begin{equation}
  F_{\bm 0,u}^{\rm ref}
  =-A_{\rm ref}
   \sum_m\overline\rho_{f,u,m}
   \sum_\nu\Cchan_{m,\nu}
   \Phiint\!\left(Q_{z,m}^{\nu};z_m^-,z_m^+\right),
  \label{eq:per-cell-reference}
\end{equation}
so that record $u$'s complete share of the specular matrix element is
\begin{equation}
  \boxed{
  F_{\bm 0,u}
  =\sum_m\sum_\nu\mathcal A_{u,m}^{\nu}
   -A_{\rm ref}
    \sum_m\overline\rho_{f,u,m}
    \sum_\nu\Cchan_{m,\nu}
    \Phiint\!\left(Q_{z,m}^{\nu};z_m^-,z_m^+\right).
  }
  \label{eq:per-cell-specular}
\end{equation}
Summing Eq.~\eqref{eq:per-cell-specular} over $u$ and applying
Eq.~\eqref{eq:reference-partition} recovers
Eq.~\eqref{eq:profile-reference-term} and hence the profile-level
$F_{\bm 0}$ of Eq.~\eqref{eq:Fh-two-parts}:
\begin{equation}
  \boxed{
  F_{\bm 0}
  =\sum_uF_{\bm 0,u}
  =\sum_uF_{\bm 0,u}^{\rm at}
   -A_{\rm ref}
    \sum_m\overline\rho_{f,m}
    \sum_\nu\Cchan_{m,\nu}
    \Phiint\!\left(Q_{z,m}^{\nu};z_m^-,z_m^+\right).
  }
  \label{eq:total-specular}
\end{equation}
If callers do not need an attribution by object,
Eq.~\eqref{eq:total-specular} is the preferred calculation: accumulate all
atomic terms, evaluate the prepared reference once, and add once.

\section{Reduction to the implemented terminal bulk formula}

Let the top terminal bulk cell occupy $b-c\leq z<b$, with repeat vector
$-\bm c$ into the substrate and $c=c_z>0$.  Only the $++$ channel survives, by
Eq.~\eqref{eq:terminal-branches}.  If $F_{\rm uc}$ denotes the ordinary
unit-cell structure factor including the global phase of the top cell, the
semi-infinite atomic contribution to $F_{\bm 0}$ is
\begin{equation}
  F_{\bm 0,\rm bulk}^{\rm at}
  =\Cchan_{M,++}
   \frac{F_{\rm uc}}
        {1-\exp(-\ii\bm Q_M^{++}\mathbin{\cdot}\bm c)},
  \label{eq:bulk-atomic-sum}
\end{equation}
and, for a surface-normal repeat, the mean-density contribution is
\begin{equation}
  F_{\bm 0,\rm bulk}^{\rm ref}
  =-\Cchan_{M,++}A_{\rm ref}\overline\rho_{f,M}
   \frac{\ee^{\ii Q_{z,M}^{++}b}}{\ii Q_{z,M}^{++}}.
  \label{eq:bulk-reference}
\end{equation}
Their sum is the previously implemented bulk result.  The four-channel
finite-layer formula therefore extends the existing kernel without changing
the terminal half-space convention.

Equation~\eqref{eq:bulk-reference} also shows why the reference contribution
cannot be dropped: it is precisely the $Q_z\to0$ divergence of the kinematic
lattice sum in Eq.~\eqref{eq:bulk-atomic-sum}, that is, the $\bm h=\bm 0$
component that the prepared optical profile already describes and that $r_0$
in Eq.~\eqref{eq:specular-reflectivity} already carries.  Retaining both
without Eq.~\eqref{eq:total-specular} would count it twice.
